\documentclass{article}
\usepackage{xcolor}
\newif\ifaddcomments\addcommentstrue
\newcommand{\todo}[1]{\ifaddcomments{\textcolor{red}{[TODO: #1]}}\fi}
\newcommand{\alice}[1]{\ifaddcomments{\textcolor{blue}{[Alice: #1]}}\fi}
\newcommand{\bob}[1]{\ifaddcomments{\textcolor{violet}{[Bob: #1]}}\fi}
\newcommand{\carol}[1]{\ifaddcomments{\textcolor{teal}{[Carol: #1]}}\fi}

\begin{document}

\section{Introduction}
Learned action models make planning possible without a hand-written domain.
We show that a policy cloned from demonstrations can guide a planner on problems that are far from the demonstrations.
\alice{``far from the demonstrations'': say how far, or point to the section that defines the distance.}
\bob{Section 5 defines it; I will add the reference.}
\alice{Fine, but also give one number in the introduction.}

\section{Experiments}
\subsection{Setup}
The policy is a two-layer MLP with 128 hidden units.
The policy is trained for 50 epochs on the demonstrations.
\carol{``trained for 50 epochs'': how was the number of epochs chosen? Validation split, early stopping?}

Each configuration gets 60 seconds per problem \todo{check this is the same for both planners}.

\subsection{Results}
Figure 2 shows success against problem size.
% \alice{Report the number of runs behind each point.}
% \bob{Done: 10 runs per point, now in the caption.}

\section{Conclusions}
First, we test only small domains.
Second, our domains are all small.
\carol{The second limitation repeats the first one.}

\end{document}
